\documentclass[11pt]{article}
\usepackage[margin=1in]{geometry}
\usepackage{hyperref}
\usepackage{enumitem}
\usepackage{longtable}
\usepackage{tabularx}

\title{AICT As-Built Notes}
\author{}
\date{2026-01-28}

\begin{document}
\maketitle
\tableofcontents

\section{Overview}
This document captures the current implementation state of the AI Canvas extension host
components added after the implementation specification. The focus is the RPC server and the
host-side subsystems that support it (policy, patching, runner, repo index, storage, cloud gateway).

\section{Status Summary}
\begin{itemize}[noitemsep]
  \item RPC server is implemented and wired in the extension activation flow.
  \item Cloud gateway uses Anthropic (Claude) only; OpenAI and Gemini adapters are stubs.
  \item Repo indexing, patch application, and policy enforcement are minimal but functional.
  \item Shared types and Zod schemas are in place under \texttt{src/shared}.
\end{itemize}

\section{Folder Layout (Implemented)}
\begin{itemize}[noitemsep]
  \item \texttt{src/extension/rpc}: RPC transport, validation, server, and handlers.
  \item \texttt{src/extension/policy}: guardrails for command allowlist, scope fence, dependency gate.
  \item \texttt{src/extension/patch}: diff validation, dry-run/apply, optional formatter runner.
  \item \texttt{src/extension/runner}: command execution and output truncation.
  \item \texttt{src/extension/repoIndex}: file scan, manifest scan, import graph, symbol table.
  \item \texttt{src/extension/storage}: workspace/cache read-write helpers.
  \item \texttt{src/extension/cloud}: Claude adapter, gateway, validators, context packager.
  \item \texttt{src/shared}: shared types and Zod schemas.
\end{itemize}

\section{RPC Server}
\subsection{Core Modules}
\begin{itemize}[noitemsep]
  \item \texttt{rpcTransport.ts}: wraps \texttt{webview.postMessage} and message listeners.
  \item \texttt{rpcValidation.ts}: validates requests/responses with Zod.
  \item \texttt{rpcServer.ts}: dispatches validated requests to handlers and emits events.
  \item \texttt{rpcHandlers.ts}: method handlers for \texttt{startWork}, \texttt{runTests}, \texttt{applyPatch},
        \texttt{repoIndex}, \texttt{exportBundle}, \texttt{cancelJob}.
\end{itemize}

\subsection{Event Flow}
\begin{itemize}[noitemsep]
  \item UI sends \texttt{RpcRequest}. The server validates then routes to the handler.
  \item Handlers emit \texttt{RpcEvent} for job status, logs, and repo index updates.
  \item Responses are sent with \texttt{RpcResponse} and structured errors on failure.
\end{itemize}

\section{Policy Engine}
\begin{itemize}[noitemsep]
  \item \texttt{commandAllowlist.ts}: pattern-based allowlist using minimatch.
  \item \texttt{scopeFence.ts}: path allow/deny checks and max file limits.
  \item \texttt{dependencyGate.ts}: detects manifest and lockfile edits in diffs.
  \item \texttt{networkPolicy.ts}: allows/blocks network usage.
  \item \texttt{policyEngine.ts}: aggregates decisions and returns allow/deny with reasons.
\end{itemize}

\section{Patch Engine}
\begin{itemize}[noitemsep]
  \item \texttt{diffValidator.ts}: validates unified diff structure and scope rules.
  \item \texttt{patchApplier.ts}: applies unified diffs using \texttt{diff} parser.
  \item \texttt{formatRunner.ts}: runs format/lint commands (allowlisted).
  \item \texttt{patchEngine.ts}: orchestrates validation, dry-run, format, and apply.
\end{itemize}

\section{Runner}
\begin{itemize}[noitemsep]
  \item \texttt{commandRunner.ts}: spawns commands with timeouts, output capture, truncation.
  \item \texttt{outputTruncation.ts}: head/tail truncation for large outputs.
  \item \texttt{runQueue.ts}: basic concurrency-limited queue (not currently wired).
\end{itemize}

\section{Repo Index}
\begin{itemize}[noitemsep]
  \item \texttt{manifestScanner.ts}: parses \texttt{package.json}, detects other manifests.
  \item \texttt{importGraph.ts}: TS/JS import graph via TypeScript compiler API.
  \item \texttt{symbolTable.ts}: exported symbol discovery.
  \item \texttt{testCommandDiscovery.ts}: derives test commands from npm scripts.
  \item \texttt{repoIndexer.ts}: orchestrates scanning and returns \texttt{RepoIndex}.
\end{itemize}

\section{Storage}
\begin{itemize}[noitemsep]
  \item \texttt{workspaceStore.ts}: loads/saves \texttt{.vibecanvas.json} with schema validation.
  \item \texttt{cacheStore.ts}: loads/saves \texttt{.vibecanvas.cache.json}.
  \item \texttt{storagePaths.ts}: path helpers for workspace/cache files.
  \item \texttt{storageTypes.ts}: workspace/cache types.
\end{itemize}

\section{Cloud Gateway (Claude Only)}
\begin{itemize}[noitemsep]
  \item \texttt{claudeAdapter.ts}: Anthropic Messages API with plan/diff validation.
  \item \texttt{gatewayClient.ts}: rate-limited wrapper and context bundling.
  \item \texttt{contextPackager.ts}: size-capped context files bundle.
  \item \texttt{outputValidators.ts}: plan JSON and unified diff validation.
  \item \texttt{openaiAdapter.ts}, \texttt{geminiAdapter.ts}: stubs that throw.
\end{itemize}

\subsection{Environment}
\begin{itemize}[noitemsep]
  \item \texttt{ANTHROPIC\_API\_KEY} must be set in the extension host environment.
\end{itemize}

\section{Extension Wiring}
The command \texttt{ai-canvas.openCanvas} creates a webview panel, then constructs:
\begin{itemize}[noitemsep]
  \item \texttt{RpcTransport} for the panel webview.
  \item \texttt{PolicyEngine}, \texttt{PatchEngine}, \texttt{CloudGateway}, \texttt{RepoIndexer}.
  \item \texttt{RpcServer} with core handlers and starts it.
\end{itemize}

\section{Dependencies}
Runtime dependencies added:
\begin{itemize}[noitemsep]
  \item \texttt{zod}, \texttt{minimatch}, \texttt{fast-glob}, \texttt{parse-diff}, \texttt{diff}.
\end{itemize}

Dev dependencies added:
\begin{itemize}[noitemsep]
  \item \texttt{@types/diff}.
\end{itemize}

\section{Build and Verification}
\begin{itemize}[noitemsep]
  \item TypeScript compile: \texttt{npm run compile-tests}
  \item Extension bundle: \texttt{npm run compile}
\end{itemize}

\end{document}
